% ============================================================
% APPENDICE B — Template Pronti all'Uso
% ============================================================
\chapter{Template Pronti all'Uso}

Tutti i template di questo libro, pronti da copiare e adattare ai tuoi progetti.

% ---------------------------------------------------------
\section{Template \texttt{\_CONTEXT.md} Base}

\begin{lstlisting}[language={},title={\_CONTEXT.md --- Template base}]
# Progetto: [Nome Progetto]

## Scopo
[Descrizione in 2-3 frasi di cosa fa l'applicazione]

## Tecnologie
- Linguaggio: [es. Python 3.12, Node.js 20, Dart]
- Framework: [es. Express.js, React, Flutter]
- Database: [es. PostgreSQL 16 con Prisma]
- Altro: [dipendenze chiave]

## Architettura
[Diagramma ASCII o descrizione della struttura]

src/
  [struttura directory con commento per ogni cartella]

## Convenzioni
- [Naming conventions]
- [Pattern architetturali]
- [Formato risposta API se applicabile]

## Vincoli
- [NON fare X perche Y]
- [SEMPRE fare Z]
- [Limiti tecnici o di sicurezza]

## Comandi
- Avviare: [comando]
- Test: [comando]
- Build: [comando]
\end{lstlisting}

% ---------------------------------------------------------
\section{Template \texttt{\_CONTEXT.md} REST API}

\begin{lstlisting}[language={},title={\_CONTEXT.md --- REST API}]
# Progetto: [Nome] API

## Scopo
[Descrizione del servizio API]

## Tecnologie
- Runtime: Node.js 20
- Framework: Express.js
- Validazione: Zod
- Database: PostgreSQL 16 con Prisma ORM
- Auth: JWT + OAuth 2.0
- Docs: Swagger/OpenAPI

## Architettura
src/
  index.js          <- Entry point, avvio server
  app.js            <- Express app, middleware, route
  routes/           <- Definizione endpoint
  controllers/      <- Gestione request/response
  services/         <- Logica business
  middleware/       <- Auth, validazione, error handling
  prisma/           <- Schema e migrazioni

## Endpoint
| Metodo | Path            | Auth | Descrizione     |
|--------|-----------------|------|-----------------|
| GET    | /api/resource   | Si   | Lista risorse   |
| POST   | /api/resource   | Si   | Crea risorsa    |
| GET    | /api/resource/:id | Si | Dettaglio       |
| PUT    | /api/resource/:id | Si | Aggiorna        |
| DELETE | /api/resource/:id | Si | Elimina         |

## Formato Risposte
Successo: { "success": true, "data": { ... } }
Errore:   { "success": false, "error": { "message": "..." } }

## Vincoli
- Validazione Zod PRIMA del controller
- Ogni query filtra per userId (autorizzazione)
- Errori in produzione: solo messaggi generici
- Rate limiting su endpoint auth

## Comandi
- Dev: npm run dev
- Test: npx vitest run
- Migrate: npx prisma migrate dev
- Studio: npx prisma studio
\end{lstlisting}

% ---------------------------------------------------------
\section{Template \texttt{SKILL.md} React}

\begin{lstlisting}[language={},title={SKILL.md --- React Developer}]
---
name: react-developer
description: Skill per sviluppo frontend React con Tailwind CSS.
---

# React Developer Skill

## Stack
- React 18+ con Vite
- Tailwind CSS 4
- React Router 6
- Axios per HTTP

## Struttura
src/
  components/ui/     <- Componenti riutilizzabili
  components/layout/ <- Header, Footer, Container
  pages/             <- Pagine (una per route)
  hooks/             <- Hook personalizzati
  context/           <- Context API per stato globale
  services/          <- Chiamate API (Axios)

## Convenzioni
- Componenti: PascalCase, functional only
- Hook: camelCase con prefisso "use"
- NON fare fetch nei componenti, usa hook
- NON usare localStorage per token JWT
- Ogni pagina: loading + error + empty state
\end{lstlisting}

% ---------------------------------------------------------
\section{Template \texttt{SKILL.md} Flutter}

\begin{lstlisting}[language={},title={SKILL.md --- Flutter Developer}]
---
name: flutter-developer
description: Skill per sviluppo app mobile Flutter.
---

# Flutter Developer Skill

## Stack
- Flutter 3.x / Dart
- Riverpod 2 per state management
- GoRouter per navigazione
- Dio per HTTP
- flutter_secure_storage per token

## Struttura
lib/
  config/      <- Configurazione (API URL, tema)
  models/      <- Classi immutabili con fromJson/toJson
  services/    <- Client HTTP, auth, business logic
  providers/   <- Riverpod StateNotifier
  screens/     <- Schermate (una per route)
  widgets/     <- Widget riutilizzabili

## Convenzioni
- File: snake_case.dart, Classi: PascalCase
- Widget: ConsumerWidget per accedere ai provider
- Modelli: classi immutabili con factory constructor
- NON usare setState per stato condiviso
- NON usare SharedPreferences per token
- Material Design 3 con useMaterial3: true
\end{lstlisting}

% ---------------------------------------------------------
\section{Template \texttt{SKILL.md} Analisi}

\begin{lstlisting}[language={},title={SKILL.md --- Discovery \& Analysis Specialist}]
---
name: analysis-specialist
description: Skill per fasi Discovery (0) e Analysis (1).
---

# Discovery & Analysis Specialist

## Principi
- Technology-agnostic: definisci COSA, mai COME
- Requisiti verificabili: ogni requisito ha un criterio di accettazione
- Sicurezza e Performance obbligatorie (SEC/PERF)
- Scope discipline: definire esplicitamente OUT OF SCOPE

## Output Attesi
- Glossario condiviso
- Requisiti funzionali (FR-01..FR-NN)
- NFR: Sicurezza (SEC-01..SEC-05) + Performance (PERF-01..PERF-05)
- User story con GIVEN/WHEN/THEN
- Modello dati concettuale (NO tipi fisici)
- Threat model preliminare

## Formato User Story
[US-XXX] Titolo
As a [RUOLO] I want [AZIONE] So that [BENEFICIO]
Acceptance Criteria:
- GIVEN [contesto] WHEN [azione] THEN [risultato]

## Vincoli
- MAI scegliere tecnologie (e' compito della fase Design)
- MAI scrivere codice o pseudocodice
- MAI saltare i requisiti NFR di sicurezza e performance
- MAI creare una user story senza almeno un criterio GIVEN/WHEN/THEN
\end{lstlisting}

% ---------------------------------------------------------
\section{Template \texttt{SKILL.md} Design}

\begin{lstlisting}[language={},title={SKILL.md --- Architecture \& Design Specialist}]
---
name: design-specialist
description: Skill per fase Design (2).
---

# Architecture & Design Specialist

## Principi
- Ogni decisione tracciabile a un requisito (FR o NFR)
- Trade-off trasparenti: documenta pro, contro, rationale in ADR
- Security by design e Performance by design
- Task piccoli e stimabili (max 8 story point)

## Output Attesi
- Valutazione stack (matrice pesata)
- ADR per ogni scelta tecnologica significativa
- Architettura di sicurezza e performance
- Contratto API (OpenAPI o equivalente)
- Schema database con indici
- EPIC con task breakdown e Definition of Done

## Formato ADR
# ADR-NNN: [Titolo]
Status: Proposed | Accepted | Deprecated
Date: YYYY-MM-DD
## Context - ## Decision - ## Consequences

## Vincoli
- MAI iniziare il design senza requisiti approvati
- MAI scegliere stack senza matrice di valutazione pesata
- MAI creare task piu' grandi di 8 story point
- SEMPRE scrivere un ADR per ogni scelta tecnologica
\end{lstlisting}

% ---------------------------------------------------------
\section{Template \texttt{SKILL.md} Operations}

\begin{lstlisting}[language={},title={SKILL.md --- Operations \& Monitoring Specialist}]
---
name: ops-specialist
description: Skill per fase Operations (6).
---

# Operations & Monitoring Specialist

## Principi
- Observe, don't guess: decisioni basate su metriche e log
- Automate repeatables: se un runbook gira piu' di 2 volte, scriptalo
- Blameless post-mortem: focus sui sistemi, non sulle persone
- Recovery over root cause: ripristina il servizio PRIMA

## SLA Target
- Availability: 99.9% (alert < 99.5%)
- Response Time P95: < 200ms (alert > 500ms)
- Error Rate: < 0.1% (alert > 1%)

## Incident Response
1. ALERT TRIGGERED
2. ACKNOWLEDGE (entro 5 min)
3. CLASSIFY SEVERITY (P1 Critical -> P4 Low)
4. INVESTIGATE (dashboard, deploy recenti, log)
5. MITIGATE (rollback, scale, feature-flag)
6. ROOT CAUSE ANALYSIS (solo dopo ripristino)
7. POST-MORTEM (blameless, action items)
8. UPDATE RUNBOOK

## Vincoli
- MAI applicare fix senza verificare il risultato
- MAI saltare il post-mortem dopo un incidente P1/P2
- MAI modificare dati di produzione senza backup
- MAI silenziare un alert senza creare un ticket
- SEMPRE ripristinare il servizio PRIMA di investigare
\end{lstlisting}

% ---------------------------------------------------------
\section{Template \texttt{\_CONTEXT.md} Full-Stack}

\begin{lstlisting}[language={},title={\_CONTEXT.md --- Full-Stack}]
# Progetto: [Nome] Full-Stack

## Panoramica
[Descrizione in 3-4 frasi]

## Architettura
[Browser] -> [Frontend :5173] -> [Proxy /api]
  -> [Backend :3000] -> [PostgreSQL]

## Componenti

### Backend (./backend/)
- Contesto: ./backend/_CONTEXT.md
- Porta: 3000

### Frontend (./frontend/)
- Contesto: ./frontend/_CONTEXT.md
- Skill: ./frontend/SKILL.md
- Porta: 5173

### Mobile (./mobile/) -- opzionale
- Contesto: ./mobile/_CONTEXT.md
- Skill: ./mobile/SKILL.md

## Contratto API
[Tabella endpoint condivisa tra tutti i componenti]

## Variabili d'Ambiente
### Backend: DATABASE_URL, JWT_SECRET, ...
### Frontend: VITE_API_URL

## Comandi
- Tutto: npm run dev (con concurrently)
- Backend: cd backend && npm run dev
- Frontend: cd frontend && npm run dev
\end{lstlisting}

% ---------------------------------------------------------
\section{Template \texttt{PROGRESS.md}}

\begin{lstlisting}[language={},title={PROGRESS.md}]
# [Nome Progetto] -- Progresso

## Stato Attuale
- [x] Task completato 1
- [x] Task completato 2
- [ ] Task in corso
- [ ] Task futuro

## Decisioni Architetturali
- [Decisione]: [Motivazione]

## Problemi Risolti
- [Problema]: [Soluzione adottata]

## Note per la Prossima Sessione
- [Cosa ricordare per continuare il lavoro]
\end{lstlisting}

% ---------------------------------------------------------
\section{Template \texttt{.gitignore} Full-Stack}

\begin{lstlisting}[language={},title={.gitignore}]
# Environment
.env
.env.local
.env.production

# Dependencies
node_modules/

# Build
dist/
build/

# Database
*.db
*.sqlite

# IDE
.vscode/settings.json
.idea/

# OS
.DS_Store
Thumbs.db

# Secrets
*.jks
*.keystore
key.properties

# Flutter
.dart_tool/
.packages
\end{lstlisting}

% ---------------------------------------------------------
\section{Checklist Pre-Deploy}

\begin{lstlisting}[language={},title={Checklist pre-deploy}]
## Sicurezza
- [ ] JWT_SECRET >= 256 bit, in variabile d'ambiente
- [ ] Token scade (access: 1h, refresh: 7d)
- [ ] Logout invalida refresh token
- [ ] Input validato con Zod
- [ ] Errori non espongono stack trace
- [ ] Helmet.js attivo
- [ ] Rate limiting su auth
- [ ] CORS solo domini consentiti
- [ ] npm audit clean
- [ ] .env nel .gitignore

## Funzionalita
- [ ] Login OAuth funziona
- [ ] CRUD completo funziona
- [ ] Protezione route funziona
- [ ] Error handling end-to-end
- [ ] Mobile connesso al backend
- [ ] Health check endpoint

## Infrastruttura
- [ ] Database migrato in produzione
- [ ] Variabili d'ambiente impostate
- [ ] HTTPS attivo
- [ ] Monitoring attivo
\end{lstlisting}
